\documentclass[11pt,a4paper]{article}
\usepackage[margin=1in]{geometry}
\usepackage{hyperref}
\usepackage{enumitem}
\usepackage{titlesec}
\usepackage{booktabs}
\usepackage{graphicx}

\hypersetup{
    colorlinks=true,
    linkcolor=blue,
    urlcolor=blue,
}

\title{\textbf{Order Supervisor POC} \\ \large Deliverables Package}
\author{Aman Anand}
\date{\today}

\begin{document}
\maketitle

\section*{Project Overview}

A long-running AI supervisor that oversees a single e-commerce order from creation to completion. An LLM agent (Llama 3.3 via Groq API) is orchestrated by a Temporal workflow, wakes up on schedule or important events, and takes actions like sending messages, creating notes, or escalating issues. A Next.js dashboard lets operators inject events, add instructions, and monitor the agent in real-time.

\section*{Deliverables}

\subsection*{1. Source Code}
The full source code is available in the GitHub repository:

\url{https://github.com/AmanAnand958/Order-Supervisor-POC}

\noindent Repository structure:
\begin{itemize}[nosep]
    \item \texttt{backend/} — Python 3.12 + FastAPI + Temporal worker + PostgreSQL
    \item \texttt{frontend/} — Next.js 16 + React 19 + Tailwind CSS 4 + TypeScript
    \item \texttt{docker-compose.yml} — Local development (Postgres + Temporal)
    \item \texttt{docker-compose.prod.yml} — Production deployment
\end{itemize}

\subsection*{2. README with Setup Instructions}
Detailed setup instructions are provided in \texttt{README.md} within the repository, covering:
\begin{itemize}[nosep]
    \item Prerequisites (Docker, Python 3.12+, Node.js 20+, Groq API key)
    \item Step-by-step local development setup (7 steps)
    \item End-to-end usage flow
    \item Key commands reference
\end{itemize}

\subsection*{3. Architecture Note}
A comprehensive architecture document is provided in \texttt{ARCHITECTURE.md}, covering:
\begin{itemize}[nosep]
    \item Temporal workflow design (signals, queries, main loop)
    \item Sleep/wake mechanism with configurable aggressiveness
    \item Deterministic event severity classifier (no LLM calls)
    \item Rolling memory compaction approach
    \item Activity architecture (6 activities)
    \item Database role and live state querying
    \item Key tradeoffs given the POC scope
\end{itemize}

\subsection*{4. Walkthrough Video}
A walkthrough video demonstrating the product is available at:

\vspace{0.5em}
\noindent\textbf{[VIDEO LINK — TO BE INSERTED]}

\vspace{0.5em}
\noindent The video demonstrates:
\begin{enumerate}[nosep]
    \item Creating a supervisor configuration
    \item Starting an order run
    \item Sending events into the workflow
    \item The agent going to sleep and waking up on important events
    \item Tool execution (send message, create note, escalate issue)
    \item Adding extra instructions to a live run
    \item Interrupting and resuming a run
    \item Final summary, learnings, and recommendations
\end{enumerate}

\section*{Tech Stack}

\begin{center}
\begin{tabular}{ll}
\toprule
\textbf{Layer} & \textbf{Technology} \\
\midrule
Orchestration & Temporal (Python SDK) \\
Backend API & Python 3.12 + FastAPI + asyncpg \\
LLM Runtime & Groq API (llama-3.3-70b-versatile) \\
Database & PostgreSQL 16 \\
Frontend & Next.js 16 + React 19 + Tailwind CSS 4 + TypeScript \\
Containerization & Docker + Docker Compose \\
\bottomrule
\end{tabular}
\end{center}

\section*{Links}

\begin{itemize}[nosep]
    \item GitHub: \url{https://github.com/AmanAnand958/Order-Supervisor-POC}
    \item Walkthrough Video: \textbf{[VIDEO LINK — TO BE INSERTED]}
\end{itemize}

\end{document}
